\section{Discussion}\label{sec:discussion}

\subsection{Target-hour weather is the day-ahead lever}
The results in Section~\ref{sec:results} make clear that the dominant lever for
day-ahead accuracy is not the recent load history but the \emph{forecasted weather
at the target hour}. Twenty-four hours out, the most recent observed load carries
comparatively little information about tomorrow's load: the diurnal and weather-driven
structure of the day to come is what matters, and that structure is captured by
calendar features (exactly known in advance) and by the forecast temperature at
$t+h$. Because each per-horizon model is trained to read the exogenous features at its
own target hour rather than rolling a near-term prediction forward, it stays accurate
deep into the horizon. The ablation in Table~\ref{tab:ablation} quantifies this:
adding target-hour calendar and weather features cuts day-ahead (h=24) MAE from
$183.0$ to $110.3$~MW, and unifying the weather source pushes it further to
$76.7$~MW. The bulk of the day-ahead error that a lags-only model exhibits is, in
effect, an absence of knowledge about tomorrow's weather, not an absence of recent
load.

\subsection{Train/serve weather-source matching}
A second, more subtle lesson concerns deployment. A model can be internally
consistent and well-validated on held-out data yet still degrade in production if the
distribution of an input shifts silently between training and serving. Here the
culprit was the weather source: training on OpenWeatherMap history while serving
Open-Meteo forecasts induced a distribution mismatch that nearly doubled live
day-ahead error, from roughly $90$ to $147$~MW (Section~\ref{sec:results}). Retraining
on the same Open-Meteo source used at serving time removed most of the gap to the
operator. No change to the model architecture, features, or hyper-parameters was
involved; the single act of \emph{matching the train and serve weather source} was
responsible. For weather-informed forecasters this is an easy mistake to make and a
costly one, and it argues for treating the exogenous-data provider as a versioned,
unified part of the pipeline rather than an interchangeable commodity.

\subsection{The horizon-mismatch pitfall}
The comparison against ISO-NE also illustrates a methodological pitfall that can
inflate reported gains. It is tempting to compare a model's short-horizon nowcast
against an operator's day-ahead forecast, but this is not an apples-to-apples
contrast: the two forecasts solve different problems at different lead times. When the
comparison is made honestly --- both forecasters evaluated at the same 24-hour-ahead
horizon, on the same Open-Meteo weather --- Beacon's edge over the ISO-NE operational
forecast is a modest but real improvement of $-3.3$~MW MAE ($\approx +3.5\%$), with a
corresponding MAPE improvement and a marginally lower $R^2$ (Table~\ref{tab:iso}).
This is a genuine result, not the inflated figure a horizon mismatch would produce,
and reporting it as such is part of the contribution.

\subsection{Operational implications}
Taken together, these findings carry a practical message. ISO-NE's day-ahead demand
forecast is an operational product that grid operators rely on for unit commitment,
economic dispatch, and reserve sizing, and matching --- let alone slightly beating ---
it is a hard bar. That a keyless, fully reproducible pipeline built entirely on free
public data sources (ISO-NE Web Services and Open-Meteo) can clear that bar at
day-ahead horizon, while remaining auditable and leakage-clean
(Section~\ref{sec:results}), suggests that the gains here come from disciplined feature
construction and honest evaluation rather than from proprietary data or heavy modeling
machinery. Small MAE reductions translate into meaningful cost and reliability gains in
wholesale electricity markets, and a transparent pipeline that achieves them lowers the
barrier to adoption and independent verification.
